\documentclass[11pt,a4paper]{article}
\usepackage{curso}
\begin{document}
\portada{Ejercicios --- Lección}{10.9}{Caso de estudio: un modelo de lenguaje AFT}
\section{Ejercicios}
\begin{enumerate}[label=\textbf{E\arabic*.}]
\item (B) ¿Por qué la primera difusión solo aprendió frecuencias de letras?
\item (B) Razón de pesos de proyección entre el modelo denso y el de $k=4$.
\item (B) ¿Qué haría ganar al gradiente natural en tiempo?
\item (B) ¿Por qué la cifra de difusión es una cota superior?
\item (I) Cuente las capas de muestreo secuenciales del modelo $k=4$ y justifique el número 8.
\item (I) Con la cabeza de 71 símbolos, ¿qué fracción de los FLOPs por token es la cabeza en el modelo denso? ¿Y en el disperso?
\item (A) Ejecute \texttt{caso\_estudio.py} con otro corpus o con $k=16$ y actualice la tabla y el registro de mapeo.
\end{enumerate}
\section{Soluciones}
\begin{solucion}
\textbf{E1.} La posición enmascarada solo recibe información de las vecinas, diluida entre 64 posiciones; los sesgos locales no crecían por el ruido del gradiente.

\textbf{E2.} $98\,304/4608\approx21$.

\textbf{E3.} Una resolución analógica rápida y barata del sistema lineal (TNGD), o menos iteraciones de CG con buen precondicionador.

\textbf{E4.} MDLM optimiza una ELBO: $-\log p\le$ ELBO; la verosimilitud real es menor o igual.

\textbf{E5.} Por bloque: q/k/v en paralelo, luego o, w1 y w2: 4; dos bloques: 8.

\textbf{E6.} Cabeza: $2\cdot64\cdot71=9088$ FLOP; denso: $9088/206\,720\approx4{,}4\,\%$; disperso: $9088/19\,328\approx47\,\%$: al dispersar el cuerpo, la cabeza pesa mucho más
(compárese con 10.8).

\textbf{E7.} Abierto.
\end{solucion}
\end{document}
